\section{Proposed Algorithms}\label{sec:proposed_algorithms}
In this chapter, we introduce a family of algorithms which we call conditional geometric resampling (CGR).
% which includes three variants named CGR~I, CGR~II-unbiased and CGR~II-biased. 
This family is designed as an improvement over the original GR, providing an estimator of $w_{t,i}^{-1}$ with better computation efficiency.  
%We next show that these algorithm does not worsen the regret bounds of the original GR.
In the following,
we write $\sigma_i$ to denote the number of arms (including $i$ itself) whose cumulative losses do not exceed $\hat{L}_{t,i}$,
that is,
$\hat{L}_{t,i}$ is the $\sigma_i$-th smallest among $\{L_{t,j}\}_{j\in[K]}$ where $\sigma_i=\sigma_j$ can happen when $\hat{L}_{t,i}=\hat{L}_{t,j}$.
For example, a current best arm $\hat{i}_t^*\in \argmin_{j}\hat{L}_{t,j}$
is an arm $i$ such that $\sigma_i=1$.
\vspace{-0.3em}
\subsection{General Idea}
Before explaining the specific procedures of the proposed algorithms,
we first give the intuition behind them.
When we run the original GR,
some possible values of $r_t'$ clearly violates the condition for termination.
For example, when the pulled arm $I_t$ is not the current best arm $\hat{i}_t^*\in\argmin_{i}\hat{L}_{t,i}$,
then the resampled perturbation must satisfy $r'_{t,I_t}\ge r'_{t,\hat{i}_t^*}$
so that the condition for termination
\begin{align}
I_t = {\arg\min}_{i \in [K]} \{\hat{L}_{t,i} - r'_{t,i} / \eta_t\}
\label{cond_termination}
\end{align}
is satisfied since $\hat{L}_{t,I_t}\ge\hat{L}_{t,\hat{i}_t^*}$.
Thus it is somewhat ``wasteful'' to sample $r^{\prime}$ satisfying $r'_{t,I_t}< r'_{t,\hat{i}_t^*}$
and it becomes sufficient to resample $r^{\prime}$ only from those satisfying $r'_{t,I_t}\ge r'_{t,\hat{i}_t^*}$.

%In this way, if we can identify an appropriate necessary condition for termination,
%we can reduce the number of resampling by generating $r^{\prime}$ from the conditional distribution given such a necessary condition.
Now let us formalize this idea.
Let $\nec_{t}$ be an arbitrary necessary condition for termination given in \eqref{cond_termination}, which may depend on $\hat{L}_t$ and $I_t$.
For example, the above discussion corresponds to $\nec_{t}=\{r'_{t,I_t}\ge r'_{t,\hat{i}_t^*}\}$.
Then, we can easily derive the following property, whose proof is shown in Section~\ref{subsection:proof_idea}.
\begin{lemma}\label{lem: general idea}
Consider resampling of $r_t'$ from $\mathcal{D}$ conditioned on
$\nec_t$ until \eqref{cond_termination} is satisfied.
Then, the number $M_t$ of resampling satisfies
\begin{align}
\E[M_t|\hat{L}_t, I_t]=\frac{\sP[\nec_t|\hat{L}_t, I_t]}{w_{t,I_t}}.\n
\end{align}
\end{lemma}
%The proof of this lemma is shown in Appendix~\ref{subsection:proof_idea}.
By this lemma we can use $M_t/\sP[\nec_t|\hat{L}_t, I_t]$ as an unbiased estimator
of $w_{t,I_t}^{-1}$.
In addition, we can reduce the number of resampling if we take $\nec_t$ such that
$\sP[\nec_t|\hat{L}_t, I_t]$ is small.

Here note that we can trivially minimize $\sP[\nec_t|\hat{L}_t, I_t]$
by taking the termination condition \eqref{cond_termination} itself as $\nec_t$.
Still, such a choice makes it difficult to compute $\sP[\nec_t|\hat{L}_t, I_t]$
and to resample $r_t'$ from the conditional distribution given $\nec_t$.
In the following algorithms, we construct necessary conditions $\nec_t$ such that
$\sP[\nec_t|\hat{L}_t, I_t]$ is small and easily computed, while
the resampling from the conditional distribution given $\nec_t$ is also efficient.

\subsection{Conditional Geometric Resampling}
Based on the above idea, we propose three variants of conditional geometric resampling
%honda: Here we should not write ``respectively'' and please check the usage.
%respectively
called CGR~I (Algorithms~\ref{alg:CGR1}), CGR~II-unbiased and CGR~II-biased (Algorithms~\ref{alg:CGR2}).
Here, CGR~II-unbiased and
%CGR~II
II-biased consider resampling
%honda
from
%on
the same conditional distribution, and thus we will discuss them together.
In the algorithm description, \textcolor{blue}{blue parts} are the ones different from the original GR.

\paragraph{CGR~I}
This version corresponds to the case where we sample $r_t'$ conditioned on $\nec_t=\qty{r_{t,I_t}'=\max_{i:\sigma_i\leq\sigma_{I_t}}r_{t,i}'}$,
that is, the event that $r_{t,I_t}'$ is the largest among the arms whose cumulative estimated losses are no worse than arm $I_t$.
With the i.i.d.~nature of the perturbations, this conditional distribution exhibits favorable symmetry property, enabling the algorithm to efficiently generate samples by simple operations.
To be more specific, the sampling from this conditional distribution can be realized by Algorithm~\ref{alg:CGR1},
which only needs the extra value-swapping operation in addition to the original GR.
This fact and other properties of CGR~I are formalized as follows.
%Let $\sigma_i$ denote the number of arms whose cumulative losses do not exceed $\hat{L}_{t,i}$, then CGR~I samples $r_t'$ from $\mathcal{D}$ conditioned on $\nec%_t=\qty{r_{t,I_t}'=\max_{\sigma_i\leq\sigma_{I_t}}r_{t,i}'}$.
%With the i.i.d. nature of the perturbations, this conditional distribution exhibits favorable symmetry property, enabling the algorithm to efficiently generate s%amples from the desired conditional distribution through simple operations. To be specific, CGR~I initially samples $r_t'$ i.i.d. from the original distribution,% and then performs a single value-swapping operation between $r_{t,I_t}^{\prime}$ and $r_{t,i^\star}^{\prime}$,
%where $i^\star = \arg\max_{\sigma_i \leq \sigma_{I_t}} r_{t,i}^{\prime}$.
%We can see that CGR~I differs from the original geometric resampling algorithm only by adding a value-swapping operation, and thus it can be implemented for any perturbation. 
% As it can be derived that $\sP[\nec_t|\hat{L}_t, I_t]=\sigma_{I_t}$, according to Lemma~\ref{lem: general idea}, we have $\widehat{w_{t,I_t}^{-1}}=M_t\sigma_{I_t}$, which is an unbiased estimator for $w_{t,I_t}^{-1}$. 
%Algorithm~\ref{alg:CGR1} shows the complete procedure for CGR~I, and it has the following properties.
% \begin{algorithm}[t]
%     \caption{Conditional Geometric Resampling \uproman{1}}
%     \label{alg:CGR1}
%     \begin{algorithmic}[1]
%         \STATE {\bfseries Input:}
%         Chosen arm $I_t$, cumulative loss $\hat{L}_t$, learning rate $\eta_t$
%         \STATE Set $m \coloneqq 0$
%         \REPEAT
%             \STATE $m \coloneqq m + 1$
%             \STATE Sample $r_t^{\prime} = (r_{t,1}^{\prime}, r_{t,2}^{\prime}, \dots, r_{t,K}^{\prime})$ i.i.d. from $\mathcal{D}$
%             \STATE \textcolor{blue}{Swap the values of $r_{t,\rmax}^{\prime}$ and $r_{t,I_t}^{\prime}$, where \\ $\rmax={\arg\max}_{i: \sigma_i\leq\sigma_{I_t}}r_{t,i}^{\prime}\,$}
%         \UNTIL{$I_t = {\arg\min}_{i \in [K]} \{\hat{L}_{t,i} - r_{t,i}^{\prime} / \eta_t\}$}
%         \STATE Set $\widehat{w_{t,I_t}^{-1}} \coloneqq m\textcolor{blue}{\sigma_{I_t}}$
%     \end{algorithmic}
% \end{algorithm}
\begin{algorithm}[t]
\caption{Conditional Geometric Resampling I}
\label{alg:CGR1}

\KwIn{Chosen arm $I_t$, cumulative loss $\hat{L}_t$, learning rate $\eta_t$}
$m \coloneqq 0$\;

\Repeat{$I_t = {\arg\min}_{i \in [K]} \{\hat{L}_{t,i} - r_{t,i}' / \eta_t\}$}{
    $m \coloneqq m + 1$\;
    Sample $r_t' = (r_{t,1}', r_{t,2}', \dots, r_{t,K}')$ i.i.d.\ from $\mathcal{D}$\;
    {\color{blue}Swap the values of $r_{t,\rmax}'$ and $r_{t,I_t}'$, where $\rmax = {\arg\max}_{i:\,\sigma_i \le \sigma_{I_t}} r_{t,i}'$}\;
}

$\widehat{w_{t,I_t}^{-1}} \coloneqq m\,{\color{blue}\sigma_{I_t}}$\;
\end{algorithm}
\begin{lemma}\label{lem: complexity of CGR1}
The sample $r_t'$ obtained by Algorithm~\ref{alg:CGR1} follows
the conditional distribution of $\mathcal{D}$ given $\nec_t=\qty{r_{t,I_t}'=\max_{\sigma_i\leq\sigma_{I_t}}r_{t,i}'}$.
%$\mathcal{D}$ conditioned on $\nec_t=\qty{r_{t,I_t}'=\max_{\sigma_i\leq\sigma_{I_t}}r_{t,i}'}$,
In addition,
\begin{align*}
\sP_{r_t' \sim \mathcal{D}}[\nec_t|\hat{L}_t,I_t]=1/\sigma_{I_t}
\end{align*}
and the number $M_t$ of resampling satisfies
    \begin{align}
        \E_{r_t'\sim \mathcal{D}|\nec_t}[M_t|\hat{L}_t]\leq\log K+1.\n
    \end{align}
\end{lemma}
The proof of this lemma is shown in Section~\ref{subsection:proof_cgr1}.
Combined with Lemma~\ref{lem: general idea}, this result suggests $\widehat{w_{t,I_t}^{-1}}=M_t\sigma_{I_t}$ given by CGR~I is an unbiased estimator for $w_{t,I_t}^{-1}$.
Moreover, the expected number of resampling per round is bounded by $\log K+1$, which is independent of $\hat{L}_t$.

\begin{remark}{\rm%
In CGR I we only use the symmetry of the distribution of $r_t\in\mathbb{R}^K$.
As a result, the unbiasedness and Lemma~\ref{lem: complexity of CGR1} are still valid even if
the components of $r_t\in\mathbb{R}^K$ are not independent
as far as the joint distribution is symmetric, though the regret bounds discussed in this thesis are for independent perturbations.
}\end{remark}


\paragraph{CGR~II}
In this version we assume that $\mathcal{D}$ is supported over $[0,\infty)$.
CGR~II corresponds to the case where we sample
$r_t'$ from $\mathcal{D}$ conditioned on
\begin{align}
\nec_t= \qty{r_{t,I_t}'=\max_{i:\sigma_i\leq\sigma_{I_t}}r_{t,i}',\,r_{t,I_t}'\geq\eta_t\underline{\hat{L}}_{t,I_t}},
\label{A_CGR2}
\end{align}
that is, we impose the extra condition $r_{t,I_t}\ge \eta_t\underline{\hat{L}}_{t,I_t}$ in addition to the condition in CGR~I.
%which leads to to the relatively higher probability of termination.
We also set the maximum number of resampling $G_t\in \mathbb{N}\cup \{\infty\}$,
where we refer to the versions $G_t=\infty$ by CGR II-unbiased and $G_t<\infty$ by CGR II-biased.
The complete procedure of CGR~II is given in Algorithm~\ref{alg:CGR2}, which is explained below.

Let $\mathcal{D}_{m}$ be the distribution of the maximum of $m$ i.i.d.~samples from $\mathcal{D}$,
whose CDF is $F^{m}(x)=(F(x))^m$.
%which has CDF $F^{m}(x)=(F(x))^m$.
In CGR~II we initially sample $r_{t,I_t}'$ conditioned on $\nec_t$,
which follows $\mathcal{D}_{\sigma_{I_t}}$ truncated over
$[\eta_t\underline{\hat{L}}_{t,I_t}, \infty)$
%$\left[\eta_t\underline{\hat{L}}_{t,I_t}, \infty\right)$, whose
with CDF given by
\begin{align}
&    F_{I_t}(x; \eta_t\underline{\hat{L}}_{t,I_t})=\frac{F^{\sigma_{I_t}}(x)-F^{\sigma_{I_t}}(\eta_t\underline{\hat{L}}_{t,I_t})}{1-F^{\sigma_{I_t}}(\eta_t\underline{\hat{L}}_{t,I_t})},\quad x\ge \eta_t\underline{\hat{L}}_{t,I_t}.
    \label{eq: CDF of I_t}
\end{align}
%within its support.
Next, we sample $\qty{r_{t,i}^{\prime} \mid \allowbreak\sigma_i \leq \sigma_{I_t}, i \neq I_t}$ i.i.d.~from $\mathcal{D}$
truncated over $\left[0, r_{t,I_t}^{\prime}\right)$, whose CDF is given by
\begin{align}
    F_{i}(x; r_{t,I_t}')=F(x)/F(r_{t,I_t}'),
\qquad\quad x\in[0,r_{t,I_t}^{\prime}].
    \label{eq: CDF of i}
\end{align}
% Since $\sP[\nec_t|\hat{L}_t, I_t]$ in this algorithm is expressed as $\qty(1-F^{\sigma_{I_t}}\qty(\nu+\eta_t\underline{\hat{L}}_{t,I_t}))/\sigma_{I_t}$, according to Lemma~\ref{lem: general idea}, we have $\widehat{w_{t,I_t}^{-1}}=M_t\sigma_{I_t}/\qty(1-F^{\sigma_{I_t}}\qty(\nu+\eta_t\underline{\hat{L}}_{t,I_t}))$ as an estimator for $w_{t,I_t}^{-1}$. 
Note that the computation of $F\qty(x)$ and $F^{-1}\qty(x)$ is necessary in this algorithm to generate
samples from \eqref{eq: CDF of I_t} and \eqref{eq: CDF of i} by the method of inverse transform sampling.
We can show that the samples generated by this procedure indeed follows the conditional distribution given $\nec_t$.
% \begin{algorithm}[t]
%     \caption{Conditional Geometric Resampling \uproman{2}}
%     \label{alg:CGR2}
%     \begin{algorithmic}[1]
%         \STATE {\bfseries Input:} 
%         Chosen arm $I_t$, cumulative loss $\hat{L}_t$,
%        learning rate $\eta_t$, maximum number of resampling $G_t \in \mathbb{N}\cup \{\infty\}$
%         \STATE Set $m \coloneqq 0$
%         \REPEAT
%             \STATE $m \coloneqq m + 1$
%             \STATE Sample $\qty{r_{t,i}^{\prime}\mid \sigma_i>\sigma_{I_t}}$ i.i.d. from $\mathcal{D}$

%             \STATE \textcolor{blue}{Sample $r_{t,I_t}^{\prime}$ from
%             $\mathcal{D}_{\sigma_{I_t}}$ truncated over
% $[\eta_t\underline{\hat{L}}_{t,I_t}, \infty)$,
% %$\left[\eta_t\underline{\hat{L}}_{t,I_t}, \infty\right)$,
% whose CDF is \eqref{eq: CDF of I_t}}
%             \STATE \textcolor{blue}{Sample $\qty{r_{t,i}^{\prime}\mid \sigma_i\leq\sigma_{I_t},i\neq I_t}$ i.i.d. from $\mathcal{D}$ truncated over $\left[0,r_{t,I_t}^{\prime}\right]$, whose CDF is \eqref{eq: CDF of i}}


% %            \STATE \textcolor{red}{Sample $r_{t,I_t}^{\prime}$ from the truncated distribution of 
% %            $\mathcal{D}_{\sigma_{I_t}}$ with support $\left[\nu+\eta_t\underline{\hat{L}}_{t,I_t}, \infty\right)$, whose CDF is \eqref{eq: CDF of I_t}}
% %            \STATE \textcolor{red}{Sample $\qty{r_{t,i}^{\prime}\mid \sigma_i\leq\sigma_{I_t},i\neq I_t}$ i.i.d. from the truncated distribution of $\mathcal{D}$ with support $\left[\nu,r_{t,I_t}^{\prime}\right]$, whose CDF is \eqref{eq: CDF of i}}
%         \UNTIL{$I_t = {\arg\min}_{i \in [K]} \{\hat{L}_{t,i} - r_{t,i}^{\prime} / \eta_t\}$} \textcolor{blue}{or $m\geq G_t$}
%         \STATE Set $\widehat{w_{t,I_t}^{-1}} \coloneqq m\textcolor{blue}{\sigma_{I_t}/\big(1-F^{\sigma_{I_t}}(\eta_t\underline{\hat{L}}_{t,I_t})\big)}$
% %        \STATE Set $\widehat{w_{t,I_t}^{-1}} \coloneqq m\textcolor{red}{\sigma_{I_t}/\qty(1-F^{\sigma_{I_t}}\qty(\eta_t\underline{\hat{L}}_{t,I_t}))}$
% %        \STATE Set $\widehat{w_{t,I_t}^{-1}} \coloneqq m\textcolor{red}{\sigma_{I_t}/\qty(1-F^{\sigma_{I_t}}\qty(\nu+\eta_t\underline{\hat{L}}_{t,I_t}))}$
%     \end{algorithmic}
% \end{algorithm}
\begin{algorithm}[t]
\caption{Conditional Geometric Resampling II}
\label{alg:CGR2}

\KwIn{Chosen arm $I_t$, cumulative loss $\hat{L}_t$, learning rate $\eta_t$, maximum number of resampling $G_t \in \mathbb{N}\cup\{\infty\}$}
$m \coloneqq 0$\;

\Repeat{$I_t = {\arg\min}_{i \in [K]} \{\hat{L}_{t,i} - r'_{t,i} / \eta_t\}$ \textcolor{blue}{or $m \ge G_t$}}{
    $m \coloneqq m + 1$\;

    Sample $\{r'_{t,i} : \sigma_i > \sigma_{I_t}\}$ i.i.d.\ from $\mathcal{D}$\;

    {\color{blue}Sample $r'_{t,I_t}$ from $\mathcal{D}_{\sigma_{I_t}}$ truncated over $[\eta_t\underline{\hat{L}}_{t,I_t}, \infty)$,
    whose CDF is \eqref{eq: CDF of I_t}}\;

    {\color{blue}Sample $\{r'_{t,i} : \sigma_i \le \sigma_{I_t},\, i \neq I_t\}$ i.i.d.\ from $\mathcal{D}$ truncated over $[0, r'_{t,I_t}]$,
    whose CDF is \eqref{eq: CDF of i}}\;
}

$\widehat{w_{t,I_t}^{-1}} \coloneqq m{\color{blue}\sigma_{I_t}\qty(1 - F^{\sigma_{I_t}}(\eta_t\underline{\hat{L}}_{t,I_t}))}$\;
\end{algorithm}

\begin{lemma}\label{lem: complexity of CGR2}
    In CGR~II, $r_t'$ follows the conditional distribution of $\mathcal{D}$ given
$\nec_t$ in \eqref{A_CGR2}.
    In addition,
    \begin{align*}
    \sP_{r_t' \sim \mathcal{D}}[\nec_t|\hat{L}_t,I_t]=\big(1-F^{\sigma_{I_t}}(\eta_t\underline{\hat{L}}_{t,I_t})\big)/\sigma_{I_t}
    \end{align*} 
    and the number $M_t$ of resampling satisfies
    \begin{align}
        \E_{r_t' \sim \mathcal{D}|\nec_t}[M_t|\hat{L}_t]\leq\log K+1.\n
    \end{align}
    Besides, if $\mathcal{D}$
%honda: D is a distribution, not a random variable. A random variable follows a distribution but a (fixed) distribution does not follow another distribution. 
is
%follows
the Fr\'{e}chet distributions with shape $2$, then $M_t$ satisfies
    \begin{align}
        \E_{r_t' \sim \mathcal{D}|\nec_t}[M_t|\hat{L}_t,I_t]\leq K\lor 4.\n
    \end{align}
\end{lemma}
The proof of this lemma is shown in Section~\ref{subsection:proof_cgr2}.
Combined with Lemma~\ref{lem: general idea},
we see that
$$\widehat{w_{t,I_t}^{-1}}=M_t\sigma_{I_t}/\big(1-F^{\sigma_{I_t}}(\eta_t\underline{\hat{L}}_{t,I_t})\big)$$
%$$\widehat{w_{t,I_t}^{-1}}=M_t\sigma_{I_t}/\qty(1-F^{\sigma_{I_t}}\qty(\eta_t\underline{\hat{L}}_{t,I_t}))$$
%obtained by CGR~II
serves as an unbiased estimator for $w_{t,I_t}^{-1}$ if $G_t=\infty$.
In addition, the last statement of this lemma shows that
the expected number of resampling
is bounded by $K\lor 4$
in the worst case of $\hat{L}_t$ and $I_t$
regardless of the choice of $G_t$
if $\mathcal{D}$ is the Fr\'{e}chet distributions with shape $2$.
We will show in the next chapter that the choice $G_t=(K\lor 4)\log t$ is sufficient to guarantee the BOBW regret bound.
%in Section~\ref{sect_regret}.

\begin{comment}
Different from CGR~I, CGR~II samples
$r_t'$ from $\mathcal{D}$ conditioned on
$$\nec_t=\qty{r_{t,I_t}'=\qty(\nu+\eta_t\underline{\hat{L}}_{t,I_t})\vee\max_{\sigma_i\leq\sigma_{I_t}}r_{t,i}'},$$
\honda{the inner $=$ should be $\ge$?}
which leads to to the relatively higher probability of termination.
To generate samples following such conditional distribution, CGR~II initially samples $\qty{r_{t,i}^{\prime} \mid \sigma_i > \sigma_{I_t}}$ i.i.d. from the original distribution $\mathcal{D}$ and $r_{t,I_t}^{\prime}$ from the truncated distribution of $\mathcal{D}_{\sigma_{I_t}}$ with support $\left[\nu+\eta_t\underline{\hat{L}}_{t,I_t}, \infty\right)$. 
Here, $\mathcal{D}_{\sigma_{I_t}}$ denotes the distribution of the maximum order statistic $\max_{\sigma_i \leq \sigma_{I_t}} r_{t,i}^{\prime}$, whose cumulative distribution function is $F^{\sigma_{I_t}}\qty(x)$. Given $\underline{\hat{L}}_{t,I_t}$, $r_{t,I_t}'$ follows the cumulative distribution function expressed as 
\begin{align}
    F_{I_t}(x; \eta_t\underline{\hat{L}}_{t,I_t})=\frac{F^{\sigma_{I_t}}(x)-F^{\sigma_{I_t}}(\nu+\eta_t\underline{\hat{L}}_{t,I_t})}{1-F^{\sigma_{I_t}}(\nu+\eta_t\underline{\hat{L}}_{t,I_t})}
    \label{eq: CDF of I_t}
\end{align}
within its support. Then, $\qty{r_{t,i}^{\prime} \mid \sigma_i \leq \sigma_{I_t}, i \neq I_t}$ are sampled i.i.d. from the truncated distribution of $\mathcal{D}$ with support $\left[\nu, r_{t,I_t}^{\prime}\right]$, whose cumulative distribution function given $r_{t,I_t}'$ is expressed as 
\begin{align}
    F_{i}(x; r_{t,I_t}')=F(x)/F(r_{t,I_t}')
    \label{eq: CDF of i}
\end{align}
within its support. 
% Since $\sP[\nec_t|\hat{L}_t, I_t]$ in this algorithm is expressed as $\qty(1-F^{\sigma_{I_t}}\qty(\nu+\eta_t\underline{\hat{L}}_{t,I_t}))/\sigma_{I_t}$, according to Lemma~\ref{lem: general idea}, we have $\widehat{w_{t,I_t}^{-1}}=M_t\sigma_{I_t}/\qty(1-F^{\sigma_{I_t}}\qty(\nu+\eta_t\underline{\hat{L}}_{t,I_t}))$ as an estimator for $w_{t,I_t}^{-1}$. 
The complete procedure for CGR~II is detailed in Algorithm~\ref{alg:CGR2}. Note that the computation of $F\qty(x)$ and $F^{-1}\qty(x)$ is necessary in the resampling process and the calculation of estimator, this algorithm is only applicable for $\mathcal{D}$ with specific forms of $F\qty(x)$ and $F^{-1}\qty(x)$. Since $\nec_t$ additionally considers the relationship between $r_{t,I_t}'$ and $\underline{\hat{L}}_{t,I_t}$, there exists an upper bound on the expected times of resampling in the worst case. 

For CGR~II, we introduce the maximum number of resampling $G_t$. If we hope to obtain an unbiased estimator for $w_{t,i}^{-1}$, we just need to consider CGR~II with $G_t=\infty$, which is called CGR~II-unbiased. If we consider the finite $G_t$, then the resampling will be forcibly terminated when the times of resampling $M_t$ reaches $G_t$, even if \eqref{cond_termination} is not satisfied. For this setting, the algorithm is called CGR~II-biased.  Although finite $G_t$ will introduce some bias to the estimator $\widehat{w_{t,I_t}^{-1}}$, we will show that it doesn't hurt the regret too much if $G_t=(K\lor 4)\log t$ and the perturbations follow the Fr\'{e}chet distribution with shape $2$. The properties of CGR~II are described as follows.

\begin{lemma}\label{lem: complexity of CGR2}
    In CGR~II, $r_t'$ follows the conditional distribution of $\mathcal{D}$ given $$\nec_t=\qty{r_{t,I_t}'=\qty(\nu+\eta_t\underline{\hat{L}}_{t,I_t})\vee\max_{\sigma_i\leq\sigma_{I_t}}r_{t,i}'}.$$ 
    In addition,
    \begin{align*}
    \sP_{r_t' \sim \mathcal{D}}[\nec_t|\hat{L}_t,I_t]=\qty(1-F^{\sigma_{I_t}}\qty(\nu+\eta_t\underline{\hat{L}}_{t,I_t}))/\sigma_{I_t}
    \end{align*} 
    and the number $M_t$ of resampling satisfies
    \begin{align}
        \E_{r_t' \sim \mathcal{D}}[M_t|\hat{L}_t]\leq\log K+1.\n
    \end{align}
    Besides, if $\mathcal{D}$
%honda: D is a distribution, not a random variable. A random variable follows a distribution but a (fixed) distribution does not follow another distribution. 
is
%follows
the Fr\'{e}chet distributions with shape $2$, then $M_t$ satisfies
    \begin{align}
        \E_{r_t' \sim \mathcal{D}}[M_t|\hat{L}_t,I_t]\leq K\lor 4.\n
    \end{align}
\end{lemma}
The proof of this lemma is shown in Appendix~\ref{subsection:proof_cgr2}. Combined with Lemma~\ref{lem: general idea}, it is shown that $$\widehat{w_{t,I_t}^{-1}}=M_t\sigma_{I_t}/\qty(1-F^{\sigma_{I_t}}\qty(\nu+\eta_t\underline{\hat{L}}_{t,I_t}))$$ provided by CGR~II serves as an unbiased estimator for $w_{t,I_t}^{-1}$ if $G_t=\infty$. In this algorithm, the expected times of resampling satisfies the same bound as that of CGR~I. Furthermore, if $\mathcal{D}$ is the Fr\'{e}chet distributions with shape $2$, the expected times given the the chosen arm is bounded by $K\lor 4$, which provides a guarantee for the worst case scenario.

\end{comment}

\subsection{Comparison between CGR I and CGR II}
Now let us consider the average complexity per round, which can be expressed as
\begin{align}
&
M_t\times
%\E[M_t|\hat{L}_t]\times
(\mbox{complexity of resampling once}).
\n
\end{align}
The first factor is analyzed depending on the algorithms and conditionings
in Lemmas~\ref{lem: complexity of CGR1} and \ref{lem: complexity of CGR2}, which are summarized
in Table~\ref{tab:bound}.
The second factor is $O(K)$ in both CGR~I and II,
but its factor usually becomes slightly larger in CGR~II due to the computation of $F$ and $F^{-1}$.
For this reason, while CGR~II always achieves smaller $M_t$ than CGR~I,
the total computational cost becomes higher in CGR~II in most cases
as shown in the experiments in Chapter~\ref{sect_exp}.
%Therefore, in empirical performance, although CGR~II always demonstrates fewer times of resampling than CGR~I, it requires a higher computational time cost, which is shown in Section~\ref{sect_exp}.

%\subsection{Regret analysis}
\section{Regret Analysis}\label{sect_regret}
%\honda{Since the algorithm part became a bit long I separated this to another section.}
In this chapter, we first show that CGR I and II-unbiased keep the regret bounds for the original GR as
a direct consequence of the unbiasedness of the estimators.
We then show that CGR II-biased also achieves a regret bound with the same order as the unbiased ones.
%for which we give a sketch of the proof.

\subsection{Regret Bounds}
As we demonstrated in the last chapter, CGR I and II-unbiased still provide the unbiased estimator for $w_{t,i}^{-1}$, and thus for $\ell_{t,i}$.
% \honda{In English writing, it is standard to type I, V, X (upper-case i, v, x) for roman 1, 5, 10. and you do not have to use $\backslash$uproman.}
We can directly apply this property to the BOBW analysis by \citet{pmlr-v201-honda23a} and \citet{pmlr-v247-lee24a}
for Fr\'{e}chet-type perturbations which we state below for completeness.
\begin{proposition}\label{prop_unbiased}
    FTPL with CGR I and II-unbiased with learning rate $\eta_t = c/\sqrt{t}$ for $c>0$ satisfies
\begin{align*}
\!\!\!
        \mathcal{R}(T) \leq \begin{cases}
            O\big(\sqrt{KT}\big)  &\text{in adversarial bandits,} \\
            O\qty(\sum_{i \ne i^*}\frac{\log T}{\Delta_i})& \text{in stochastic bandits,}
        \end{cases}
% \label{w_variance}
\end{align*}
if the perturbation distribution $\mathcal{D}$ is Fr\'echet-type and satisfies the conditions in \citet[Theorem 2]{pmlr-v247-lee24a}.
\end{proposition}
% \honda{I added an explanation in Introduction that we assume that the optimal arm is unique in the stochastic setting.}
%Therefore, the regret bound for CGR I and CGR II-unbiased can be directly obtained by the previous results~\citep{pmlr-v201-honda23a, pmlr-v247-lee24a}, where the BOBW result can be achieved by using Fr\'{e}chet-type perturbations.
This proposition is straightforward and we omit the proof.
This is because the proofs
of \citet{pmlr-v201-honda23a} and \citet{pmlr-v247-lee24a}
only use the variance bound given in \eqref{bound_variance} below for the property of the unbiased loss estimator
as far as the loss estimator
$\hat{\ell}_t$ is nonnegative,
where the nonnegativity trivially holds for CGR I and II-unbiased.


\begin{remark}\label{remark_variance}{\rm%
To be more precise, FTPL with CGR I and II-unbiased can attain a slightly better regret guarantee than the one with the original GR.
This is because the variance of the estimator $\widehat{w_{t,I_t}^{-1}}$ becomes
\begin{align}
\mathrm{Var}\qty[\widehat{w_{t,I_t}^{-1}}\middle|\hat{L}_t, I_t]
% \nn
=
\begin{cases}
\frac{1}{w_{t,I_t}^2}-\frac{1}{w_{t,I_t}}&\mbox{(original GR)},\\
\frac{1}{w_{t,I_t}^2}-\frac{1}{\sP(\nec_t)w_{t,I_t}}&\mbox{(CGR I, II-unbiased)},
\end{cases}
\label{bound_variance}
\end{align}
which is smaller in CGR I and II-unbiased.
Though the improvement is hidden in the big-O notation,
it explains the better empirical performance
of CGR given in Chapter~\ref{sect_exp}.
In particular, CGR II slightly improves the factor of the dominant term of the adversarial regret,
which is roughly described below and formalized in Section~\ref{app: CGR II}.
\begin{remark}[informal]\label{thm: informal}
CGR II has an advesarial regret bound better by $0.4 \sqrt{KT}$ than the existing bound for GR.
\end{remark}
%To be more precise, FTPL with CGR I and II-unbiased can attain a slightly better regret guarantee than the one with the original GR
%if we examine the terms hidden in the big-O notation.
%This is because the variance of the estimator $\widehat{w_{t,I_t}^{-1}}$ becomes
%\begin{align}
%\lefteqn{
%\mathrm{Var}[\widehat{w_{t,I_t}^{-1}}|\hat{L}_t, I_t]
%}
%\nn
%&=
%\begin{cases}
%\frac{1}{w_{t,I_t}^2}-\frac{1}{w_{t,I_t}}&\mbox{(original GR)},\\
%\frac{1}{w_{t,I_t}^2}-\frac{1}{\sP(\nec_t)w_{t,I_t}}&\mbox{(CGR I, II-unbiased)},
%\end{cases}\label{bound_variance}
%\end{align}
%which is smaller in the CGR I and II-unbiased.
%Though this difference does not affect the dominant terms of the regret bounds, it explains the better empirical performance
%of the proposed algorithms given in Section~\ref{sect_exp}.
}\end{remark}

In CGR II-biased, where the maximum number of resampling $G_{t}$ is finite, the estimator for $w_{t,i}^{-1}$ becomes biased.
Here, we provide the regret analysis of CGR II-biased and show that it can also achieve the BOBW result,
as stated in the following theorem.
\begin{theorem}\label{thm_biased}
    FTPL with CGR II-biased with learning rate $\eta_t = c/\sqrt{t}$ for $c>0$ and maximum number of resampling $G_{t}=(K\lor 4)\log t$ satisfies that
    \begin{equation*}
        \mathcal{R}(T) \leq \begin{cases}
            O\big(\sqrt{KT}\big)  &\text{in adversarial bandits,} \\
            O\qty(\sum_{i \ne i^*}\frac{\log T}{\Delta_i})& \text{in stochastic bandits,}
        \end{cases}
    \end{equation*}
    if
%the perturbation distribution
$\gD$ is the Fr\'{e}chet distribution with shape $2$.
\end{theorem}
% \honda{Please specify the condition on the perturbation distribution.}
This result shows that FTPL with CGR II-biased reduces both the worst-case and average complexities while preserving the same order of the regret bound when using the Fr\'{e}chet perturbation with shape $2$.
Here note that the GR by \citet{JMLR:v17:15-091} achieves $O(\sqrt{KT\log K})$ regret for the adversarial setting
with the maximum number of resampling $G_t=O(\sqrt{KT})$.
% \footnote{While $G_t = O(\sqrt{KT})$ was used for all $t$ in original GR for combinatorial bandits, it suffices to set $\Theta(\sqrt{Kt/\log K})$ for MAB problem.}.
Therefore our result shows that the BOBW guarantee is achievable by significantly better
maximum number of resampling in most cases.

\begin{remark}{\rm%
Our choice $G_t=O(K\log t)$ is slightly worse than $G_t=O(\sqrt{KT})$ in \citet{JMLR:v17:15-091}
if $T \approx K$,
but it is just because their goal is to bound the bias by $O(\sqrt{KT})$ rather than $O(\log T)$.
In fact, by following the same argument as Theorem \ref{thm_biased} we can see that
$G_t=O(K\log (T/K))$ is enough to achieve $O(\sqrt{KT})$ bias, which is no worse than $G_t=O(\sqrt{KT})$ whenever $K\le T$.
}\end{remark}



%In particular, compared with the GR by \citet{JMLR:v17:15-091} to achieve $O(\sqrt{KT\log K})$ regret for the adversarial setting,
%this result not only shows that the BOBW guarantee is achievable but also significantly reduces
%the maximum number of resampling from \textcolor{blue}{$O(\sqrt{Kt/\log K})$}\footnote{While $G_t = \Theta(\sqrt{KT/\log K})$ was used for all $t$ in original GR for combinatorial bandits, it suffices to set $\Theta(\sqrt{Kt/\log K})$ for MAB problem.}.
%\honda{Why is this different from Table 1?}


%\begin{remark}{\rm%
%%\honda{todo by Honda: This remark is likely to be deleted under the current introduction.}
%From the theoretical viewpoint, the required number of iterations of Tsallis-INF
%for BOBW guarantee is unknown though $O(1)$ iterations is empirically sufficient.
%In contrast, our results provide a formal guarantee on the worst-case complexity while achieving the BOBW regret bound. 
%Moreover, the practical complexity of our approach is much smaller than the worst-case guarantee as shown in Section~\ref{sect_exp},
%which demonstrates both the theoretical robustness and practical efficiency.
%}\end{remark}

% In contrast, our results provide a formal guarantee on the worst-case complexity for achieving the BOBW regret bound.
% On the other hand, our results show BOBW regret bound with a formal guarantee on the worst-case complexity,
% and the practical complexity is much smaller than the worst-case guarantee.

% \honda{It seems to be better to mention that: from the theoretical viewpoint the required number of iterations of Tsallis-INF
% for BOBW guarantee
% is unknown though $O(1)$ iterations is empirically sufficient.
% On the other hand, our results show BOBW regret bound with a formal guarantee on the worst-case complexity,
% and the practical complexity is much smaller than the worst-case guarantee.}


\subsection{Proof Sketch for Biased Estimator}
In this section we give a sketch of a proof of Theorem \ref{thm_biased}.
The complete proof is given in Section~\ref{append_regret}.
%To prove Theorem \ref{thm_biased},
We begin with the following regret decomposition
similar to \citet[Lemma 5]{JMLR:v17:15-091}.
% \honda{please fill this.}
\begin{lemma}\label{lem: regret decomposition}
%Define $\nec_{t,i}=\{I_t=i,\, \nec_t\}$. Then
The expected regret of FTPL with CGR II satisfies
    \begin{equation*}
        \mathcal{R}(T) \leq \sum_{t=1}^T \mathbb{E}\qty[\inp{\hat{\ell}_t}{w_t - e_{i^*}}] 
        + \sum_{t=1}^T \sum_{i\in [K]} \E\qty[w_{t,i} \qty(1-\frac{w_{t,i}}{\sP(\gA_{t}|\hat{L}_t, I_t=i)})^{G_t}].
    \end{equation*}
\end{lemma}
% \honda{To Jongyeong: please modify the descriptions around $\gA_{t,i}$ accordingly.
% It is also fine to declare somewhere in the appendix that we omit the conditioning.
% It may be also convenient to define $\nec_{t,i}=\{I_t=i,\, \nec_t\}$.}
The restriction on the maximum number of resamplings always reduces the expected value of $\widehat{w_{t,i}^{-1}}=M_t/\sP(\gA_{t}|\cdot)$ compared to the case $G_t=\infty$.
By using this fact we can show that the first term on the RHS can be directly bounded by the one
with the unbiased estimator
%without the maximum number of resampling
in \citet{pmlr-v201-honda23a}.
Therefore, it becomes sufficient to consider the second term, which is bounded as follows.

% Since the restriction on $M_t$ reduces the expected value of $\widehat{w_{t,i}^{-1}}=M_t/\sP(\gA_{t,i})$, compared to the case $G_t=\infty$, the first term on the right-hand side can be directly bounded by using results from the analysis without maximum resampling~\citep{pmlr-v201-honda23a}.
% \honda{While I don't doubt the final result, this argument is not clear because $w_{t,i}-e_{i^*}$ is negative for $i\neq i^*$.
% Then, we should rather say, e.g., ``The restriction on the maximum number of resampling always reduces the expected value of $\widehat{w_{t,i}^{-1}}=M_t/\sP(\gA_{t,i})$ compared to the case $G_t=\infty$.
% By using this fact we can show that the first term on the right-hand side can be directly bounded by the one without the maximum number of resampling
% in \citet{pmlr-v201-honda23a}.'' It is better if we have more concrete explanaation.}

%which becomes non-zero when $G_t$ is finite.
\begin{lemma}\label{lem: additive term by CGR biased}
    When $\mathcal{D}$ is the Fr\'{e}chet distributions with shape $2$ and $G_t=(K\lor 4)\log t$, FTPL with CGR II-biased satisfies
    \begin{equation*}
        \sum_{t=1}^T \sum_{i\in [K]}\E\qty[w_{t,i} \qty(1-\frac{w_{t,i}}{\sP(\gA_{t}|\hat{L}_t, I_t=i)})^{G_t}] \leq \log T.
    \end{equation*}
\end{lemma}
This lemma shows that the additional regret introduced by the bias
is at most logarithmic in $T$ independent of $K$.
%affects the regret only by an additive logarithmic term.

The key to this proof is bounding $w_{t,i}/\sP(\gA_{t}|\cdot)$ from below by a constant independent of $\hat{L}_t$.
When we consider the Fr\'echet perturbation we can derive simple bounds on $w_{t,i}$ and $\sP(\gA_{t}|\cdot)$ in terms of
$\sigma_{I_t}$ and $\eta_t\underline{\hat{L}}_{t,I_t}$.
From these bounds we can show
$w_{t,i}/\sP(\gA_{t}|\cdot)\ge 1/(K\lor 4)$ by considering the worst case of
$\sigma_{I_t}$ and $\eta_t\underline{\hat{L}}_{t,I_t}$.
We expect that this result can be generalized to more general Fr\'echet-type distribution by
using bounds on $w_t$ for various cases in \citet{pmlr-v247-lee24a}.
Still, it requires heavy case-by-case analysis
and we leave it to future work.
% \honda{To Jongyeong: please confirm the descriptions around here.}


%\textcolor{blue}{The key to this proof is bounding $w_{t,i}/\sP(\gA_{t,i})$ from below by constants independent of $t$.
%Specifically, we leverage the exact formulation of the Fr\'{e}chet distribution to derive a lower bound on $w_{t,i}$ that matches the form of the upper bound for $\sP(\gA_{t,i})$.
%While bounding $\sP(\gA_{t,i})$ from above can be easily extended to general distributions, our current analysis for the lower bound on $w_{t,i}$ for all $t\in \sN,i\in [K]$ heavily relies on the specific formulation of Fr\'{e}chet distribution.
%Nevertheless, we expect it can be generalized to Fr\'{e}chet-type distributions with index $2$ considered in \citet[Lemma 10]{pmlr-v247-lee24a}.}

